\documentclass[10pt,a4paper]{amsart}
\usepackage[margin=19mm]{geometry}
\usepackage{amsmath,amssymb,mathtools}
\usepackage[scr=rsfs,cal=euler]{mathalfa}
\usepackage{xfrac}
\usepackage[unicode,hidelinks,bookmarks=false]{hyperref}
\newcommand{\ie}{i.e.\ }
\newcommand{\cf}{cf.\ }
\newcommand{\resp}{respectively\ }
\newcommand{\Subsection}{\subsection}
\providecommand{\nobreakdash}{\nobreak\hbox{-}}
\setlength{\emergencystretch}{2em}
\multlinegap0pt
\usepackage{bm}
\usepackage{bbm}
\usepackage{dsfont}
\usepackage{xspace}
\usepackage{cancel}
\usepackage{mathtools}
\usepackage{cleveref}
\usepackage{amsthm}
\makeatletter
\@ifundefined{theorem}{\newtheorem{theorem}{Theorem}}{}
\@ifundefined{lemma}{\newtheorem{lemma}[theorem]{Lemma}}{}
\@ifundefined{proposition}{\newtheorem{proposition}[theorem]{Proposition}}{}
\makeatother
\DeclareMathOperator{\diam}{diam}
\DeclareMathOperator{\img}{im}
\DeclareMathOperator{\stab}{stab}
\newcommand{\dfn}{\mathrel{\mathop:}=}
\title[Geometric models and asymptotic dimension for infinite-type ]{Geometric models and asymptotic dimension for infinite-type surface mapping class groups}
\author{Michael C. Kopreski, George Shaji}
\date{Source: arXiv:2508.06679v1 (CC BY 4.0)}
\begin{document}
\maketitle

\noindent\emph{Source and licence.} Geometric models and asymptotic dimension for infinite-type surface mapping class groups. Version arXiv:2508.06679v1 (CC BY 4.0), distributed under the Creative Commons Attribution 4.0 International licence, as recorded in the arXiv licence metadata for this exact version and shown on the abstract page https://arxiv.org/abs/2508.06679v1. Full rights notice: licenses/arxiv-2508.06679-CC-BY-4.0.txt.\par\smallskip

\noindent\textit{Editorial scope.} Counted unit: the Proposition whose source label is \texttt{prop:ccar} in the article (source0.tex lines 1475--1483, source label \texttt{prop:ccar}), delivered together with the article's own complete original proof of it (source0.tex lines 1485--1524, one \texttt{proof} environment of 40 lines). No mathematics is written, repaired or re-derived: statement and proof are the article's own text line for line. Required context retained uncounted: lem:cbaction (lines 583--587); lem:expanding\_orbit\_map (lines 648--654). Every definition, display and label of those lines is retained unchanged. Omitted from this package and not redistributed: 1--582, 588--647, 655--1474, 1484--1484, 1525--2333. Nothing inside a retained excerpt is omitted or altered: every retained body is delivered exactly as the article's lines are written, so no omission receipt of any kind accompanies this package. Citations: the retained excerpts carry no citation command at all, so no bibliography entry of the article is transcribed and no reference is redistributed. Reference adaptation: none was needed and none was made; every reference inside the delivered unit resolves inside it, so no unresolved reference or citation is produced. Printed slips kept literal: no printed word, symbol, hypothesis or number of the counted statement or of its proof was altered. Layout and class: the article's own class is \texttt{amsart} with packages including amsaddr, graphicx, ulem, graphics, xifthen, enumitem, amsmath, amssymb, amsthm, tikz-cd, braket, mathtools, mathrsfs, url; the wrapper is the lane's pinned \texttt{amsart} editorial wrapper at 10pt on A4, which restates the theorem-like environments the retained text needs and adds the article's own macro definitions that the retained text needs (\texttt{\textbackslash dfn}, \texttt{\textbackslash diam}, \texttt{\textbackslash img}, \texttt{\textbackslash stab}), with extra wrapper packages (bm, bbm, dsfont, xspace, cancel, mathtools, cleveref). The delivered numbering is the unit's own. Assets: no retained span contains a figure environment, a graphics-inclusion command, a PDF-page inclusion, a table or a TikZ picture; no image file is copied into this package, the package declares no asset dependency and no image file is redistributed with it. Licence: version arXiv:2508.06679v1 of this article is distributed under the Creative Commons Attribution 4.0 International licence, as recorded in the arXiv licence metadata for this exact version and shown on the abstract page https://arxiv.org/abs/2508.06679v1. A full rights notice is delivered at licenses/arxiv-2508.06679-CC-BY-4.0.txt. Characters: every retained excerpt is pure ASCII, so the delivered engine, XeLaTeX with its default Latin Modern text font, reports no missing character.
\begin{lemma}\label{lem:cbaction}
  Suppose that a topological group $G$ acts continuously and isometrically 
  on a metric space $(X,d)$. Then the orbit map $\omega$ is
  bornologous.
\end{lemma}

\begin{lemma}\label{lem:expanding_orbit_map}
    Suppose that a topological group $G$
    acts continuously and isometrically on a metric space $(X,d)$
    and let $\omega$ be the orbit map based at $x_0 \in X$.
    Then $\omega$ is expanding if $A_\alpha \dfn \omega^{-1}(B_\alpha(x_0))$
    is coarsely bounded for all $\alpha > 0$. 
\end{lemma}

\begin{proposition}\label{prop:ccar}
 Let $G$ be a Polish group. If $G$ 
 admits a coarse Cayley--Abels--Rosendal graph then it is locally
 bounded; moreover, for
 any such graph $\Gamma$ the orbit map to 
 $V(\Gamma)$, with the induced metric,
 is a coarse equivalence.  
 If $G$ is non-Archimedean then the converse holds.
\end{proposition}

\begin{proof}
  Let $\omega : G \to V(\Gamma)$ denote the vertex orbit map; since
  the action is continuous, $\omega$ is continuous.  
  Moreover, since $V(\Gamma)$ is discrete, the stabilizer of a vertex in $\Gamma$ is thus a coarsely
  %\George{Isn't the stabilizer of a single vertex open only in the case of a simplicial graph?}
  bounded neighborhood of identity and $G$ is locally bounded.  We
  must show that $\omega$ is bornologous, expanding, and cobounded,
  hence a coarse equivalence: 
  the first follows from \cref{lem:cbaction} and the last from
  vertex transitivity, hence we need only check that $\omega$ is
  expanding. 

  Let $d$ denote the metric on $\Gamma$ and fix a vertex $x \in
  V(\Gamma)$; we assume $\omega$ is the orbit map based at $x$. 
  By \cref{lem:expanding_orbit_map}, 
  it suffices to show that $A_\alpha = \omega^{-1}(B_\alpha(x))$ is coarsely
  bounded.     Fix a connected bounded subgraph $\Lambda_\alpha
  \subset \Gamma$ containing $B_\alpha(x) \cap \img \omega =
  B_{\alpha}(x) \cap V(\Gamma)$ and let $A'_\alpha =
  \omega^{-1}(\Lambda_\alpha) = \{g \in G : gx \in
  V(\Lambda_\alpha)\}$.  Since $V(\Gamma)$ is discrete and $G$ acts
  vertex-transitively, the infimum of edge lengths in $\Gamma$ is
  non-zero. Hence by bounded-cocompactness, $\Lambda_\alpha / G$ is
  a finite graph, 
  %\George{Since $\Gamma$ is a metric graph with the metric topology, compactness doesn't imply finiteness. As an example, consider the countable rose where the edge lengths are $1/n$}, or
  or equivalently $\Lambda_\alpha$ intersects finitely many $G$-orbits of
  edges: the midpoints of edges in $\Lambda_\alpha/G$ are discrete, hence
  must be finite by compactness.   Let $\nu_x \dfn \stab_G(x) \leq G$ denote the
  stabilizer of $x$ and fix a finite set of elements $F_\alpha
  \subset G$ so
  that if $(gx,hx) \in E(\Lambda_\alpha)$, then $g^{-1}h \in
  \nu_x F_\alpha \nu_x$;
  additionally add some element $g_0 \in A'_\alpha$.  Let $m = \diam
  \Lambda_\alpha$.  Then $A'_\alpha \subset (F_\alpha \nu_x)^m$,
  hence $A'_\alpha \supset A_\alpha$ 
  is coarsely bounded since likewise is $\nu_x$.

  The converse when $G$ has small subgroups is shown in the
  following lemma.
\end{proof}

\end{document}
